% ─────────────────────────────────────────────────────────────────────────────
% Appendix: GNN Compatibility Prototype — Code Listings
% ─────────────────────────────────────────────────────────────────────────────
\section{GNN Compatibility Prototype: Code Listings}
\label{app:gnn-prototype}
This appendix contains the code listings for the GNN compatibility prototype
described in Section~\ref{sec:method-gnn-prototype}. The pipeline consists of
three stages: per-case \texttt{HeteroData} construction
(Listing~\ref{lst:heterodata-code}), structural validation
(Listing~\ref{lst:validation-code}), and the GraphSAGE forward pass
(Listing~\ref{lst:forwardpass-code}).

The pipeline is implemented in
\texttt{network\_singlecase\_semantic\_v7-2-3\_gnn-2\_canonical\_corpus.ipynb},
using Python~3.9, PyTorch~2.x, and PyTorch Geometric~2.x. The corpus is built
by looping the same canonical per-case construction logic used for the
single-case build in Chapter~4; no pre-serialised \texttt{.pt} bundle is
loaded. Running the canonical corpus builder on the FDO dataset produces the
following output:

\begin{verbatim}
=== Design case (4200-73111-00001-22, predictive variant) ===
Nodes: 5   Edges: 27   Distinct edge types: 21
Validation: P1 / P4 / P5 / P9: all pass

Feature dimensions (corpus-wide):
  offender : 39  (19 rf_* keys x 2 + 1 age)
  victim   : 19  ( 9 rf_* keys x 2 + 1 age)

HeteroData after T.ToUndirected():
  Node types : ['case', 'child', 'offender', 'victim']
  Edge types : 48
  Forward pass: OK  risk score = 0.0013  (untrained, 65,297 parameters)

=== Canonical corpus (28 FDO cases, predictive variant) ===
Canonical predictive HeteroData builds: 28/28 succeeded
Directed predictive nodes: 3-9 (mean=5.4)
Directed predictive edges: 21-41 (mean=29.0)
Canonical corpus forward passes: 28/28 succeeded
\end{verbatim}

\noindent All 28 cases pass structural validation and the GNN forward pass
without modification. The four survivor cases from the supplementary dataset
are processed separately and are not included in this corpus run.

\noindent\textit{Note: \texttt{child} nodes are present only in cases with
recorded children; the node-type list reflects all types instantiated across
the corpus rather than any single case. Edge counts in the directed corpus
statistics are pre-\texttt{ToUndirected}; the design-case summary of 48 edge
types is post-\texttt{ToUndirected}. The design-case directed edge count of
27 matches Table~3.}

% ─────────────────────────────────────────────────────────────────────────────
\subsection*{Listing 1: Per-Case \texttt{HeteroData} Construction}
\texttt{heterodata\_from\_canonical\_graph()} converts the canonical
\texttt{nodes\_df}/\texttt{edges\_df} output of the per-case graph builder
into a \texttt{HeteroData} object. Person nodes are split into typed slots by
role (\texttt{offender}, \texttt{victim}, \texttt{child}); all other node
types (\texttt{case}) retain their original label. Feature vector dimensions
are consistent across the corpus because all cases are processed against the
same fixed \texttt{FACTOR\_AS\_ATTR\_COLS} column list (P8). The
\texttt{label} parameter is not set in the current corpus; outcome labels must
be attached from the FDO's outcome columns before any training run, to prevent
target leakage (P4). \texttt{T.ToUndirected()} is applied inside this
function so that the \texttt{case} node receives incoming messages during the
GNN forward pass.

\begin{lstlisting}[
    language=Python,
    caption={Per-case \texttt{HeteroData} construction
             (\texttt{network\_singlecase\_semantic\_v7-2-3\_gnn-2\_canonical\_corpus.ipynb}).},
    label={lst:heterodata-code},
    basicstyle=\ttfamily\footnotesize,
    breaklines=true,
    frame=single,
    numbers=left,
    numberstyle=\tiny,
    keywordstyle=\bfseries,
    commentstyle=\itshape\color{gray}
]
def heterodata_from_canonical_graph(nodes_df, edges_df, context):
    """
    Convert canonical nodes/edges DataFrames to HeteroData.
    Person nodes are split by role: offender / victim / child.
    """
    node_records  = nodes_df.set_index("node_id").to_dict("index")
    offender_id   = context["offender_id"]
    victim_id     = context["victim_id"]
    case_node_id  = context["case_node_id"]
    child_nodes   = context["child_nodes"]

    # Node-type mapping: person roles -> typed integer indices
    node_type_map = {
        offender_id:  ("offender", 0),
        victim_id:    ("victim",   0),
        case_node_id: ("case",     0),
    }
    for i, ch_id in enumerate(child_nodes):
        node_type_map[ch_id] = ("child", i)

    # Feature vectors: age normalised [0,1], then rf_* pairs (value, observed)
    # offender dim = 39  (19 rf_* keys * 2 + 1 age)
    # victim dim   = 19  ( 9 rf_* keys * 2 + 1 age)
    hd = HeteroData()
    hd["offender"].x = torch.tensor(
        [node_feature_vector_from_node(node_records[offender_id])],
        dtype=torch.float)  # [1, 39]
    hd["victim"].x = torch.tensor(
        [node_feature_vector_from_node(node_records[victim_id])],
        dtype=torch.float)  # [1, 19]
    hd["case"].x = torch.tensor([[1.0]], dtype=torch.float)   # [1,  1]
    if child_nodes:
        hd["child"].x = torch.tensor(
            [[1.0]] * len(child_nodes), dtype=torch.float)    # [N,  1]

    # Edge index: one relation per (src_type, edge_type, dst_type) triple
    edge_lists = defaultdict(lambda: ([], []))
    for _, e in edges_df.iterrows():
        src_id, dst_id = e["source_id"], e["target_id"]
        if src_id not in node_type_map or dst_id not in node_type_map:
            continue
        src_type, src_idx = node_type_map[src_id]
        dst_type, dst_idx = node_type_map[dst_id]
        rel = str(e["edge_type"]).lower()     # e.g. "prior_violence"
        edge_lists[(src_type, rel, dst_type)][0].append(src_idx)
        edge_lists[(src_type, rel, dst_type)][1].append(dst_idx)

    for (src_type, rel, dst_type), (srcs, dsts) in edge_lists.items():
        hd[(src_type, rel, dst_type)].edge_index = torch.tensor(
            [srcs, dsts], dtype=torch.long)

    # Reverse edges so case node receives messages during forward pass
    hd = T.ToUndirected()(hd)
    return hd
\end{lstlisting}

% ─────────────────────────────────────────────────────────────────────────────
\subsection*{Listing 2: Structural Validation}
Before a case graph is added to the corpus, \texttt{validate\_canonical\_predictive()}
checks that three design principles hold: no incident nodes (P1), no
\textsc{killed} edge (P4), and no during- or post-stage edges in the
predictive variant (P5). All 28 cases passed without raising an exception.

\begin{lstlisting}[
    language=Python,
    caption={Structural validation per case
             (\texttt{network\_singlecase\_semantic\_v7-2-3\_gnn-2\_canonical\_corpus.ipynb}).},
    label={lst:validation-code},
    basicstyle=\ttfamily\footnotesize,
    breaklines=true,
    frame=single,
    numbers=left,
    numberstyle=\tiny,
    keywordstyle=\bfseries,
    commentstyle=\itshape\color{gray}
]
def validate_canonical_predictive(nodes_df, edges_df):
    """Schema checks against design principles P1, P4, P5."""
    if (nodes_df["node_type"] == "incident").any():
        raise ValueError("P1 violation: incident node found")
    if "KILLED" in set(edges_df["edge_type"]):
        raise ValueError("P4 violation: KILLED edge found")
    if "stage" in edges_df.columns:
        leak = edges_df[edges_df["stage"].isin(["during", "post"])]
        if len(leak) > 0:
            raise ValueError(
                "P5 violation: predictive graph contains during/post edges: "
                + ", ".join(sorted(leak["edge_type"].astype(str).unique()))
            )
\end{lstlisting}

% ─────────────────────────────────────────────────────────────────────────────
\subsection*{Listing 3: GraphSAGE Forward Pass}
Listing~\ref{lst:forwardpass-code} shows the model definition and canonical
corpus-level forward pass. The architecture is the same two-class design
as in the single-case build: \texttt{GNN} is a two-layer
\texttt{SAGEConv} module written against the homogeneous PyG API
\texttt{(x,~edge\_index)}; \texttt{to\_hetero()} expands its weight matrices
to one copy per edge type (\texttt{aggr=``sum''}). \texttt{RiskClassifier}
wraps the expanded GNN and appends a single shared \texttt{nn.Linear} readout
head that maps the offender node's 16-dimensional embedding to a scalar risk
score via sigmoid. Keeping the readout \emph{outside} \texttt{to\_hetero()}
is necessary: if it were placed inside \texttt{GNN}, \texttt{to\_hetero()}
would replicate it per edge type, which is not the intended behaviour.

For the corpus run, \texttt{T.ToUndirected()} is applied per case immediately
before the forward pass (not pre-applied to a saved bundle), so that the
directed corpus graphs produced by Listing~1 are used as the canonical
representation. A fresh \texttt{RiskClassifier} is instantiated per case
because edge-type sets can differ across cases before training-time
standardisation. The forward pass is run in inference mode with randomly
initialised weights across all 28 FDO cases; its purpose is to verify that
every \texttt{HeteroData} object in the corpus is a valid input to the
\texttt{to\_hetero()} framework.

\begin{lstlisting}[
    language=Python,
    caption={GraphSAGE forward pass
             (\texttt{network\_singlecase\_semantic\_v7-2-3\_gnn-2\_canonical\_corpus.ipynb}).},
    label={lst:forwardpass-code},
    basicstyle=\ttfamily\footnotesize,
    breaklines=true,
    frame=single,
    numbers=left,
    numberstyle=\tiny,
    keywordstyle=\bfseries,
    commentstyle=\itshape\color{gray}
]
import torch
import torch.nn.functional as F
import torch_geometric.transforms as T
from torch_geometric.nn import SAGEConv, to_hetero

HIDDEN_DIM = 16   # embedding dimension for each node type

class GNN(torch.nn.Module):
    """Two-layer GraphSAGE base (homogeneous API).
    to_hetero() expands conv1/conv2 to one weight copy per edge type."""
    def __init__(self, hidden_dim):
        super().__init__()
        self.conv1 = SAGEConv((-1, -1), hidden_dim)
        self.conv2 = SAGEConv((-1, -1), hidden_dim)

    def forward(self, x, edge_index):          # homogeneous inputs;
        x = F.relu(self.conv1(x, edge_index))  # to_hetero dispatches
        x = self.conv2(x, edge_index)          # per edge type
        return x                               # embeddings per node type

class RiskClassifier(torch.nn.Module):
    """Wraps GNN with to_hetero() and adds a per-offender readout head.
    The Linear classifier is kept outside to_hetero() so it is shared
    across all edge types rather than replicated."""
    def __init__(self, hidden_dim, metadata):
        super().__init__()
        self.gnn        = to_hetero(GNN(hidden_dim), metadata, aggr="sum")
        self.classifier = torch.nn.Linear(hidden_dim, 1)

    def forward(self, x_dict, edge_index_dict):
        embeddings   = self.gnn(x_dict, edge_index_dict)
        offender_emb = embeddings["offender"]   # [1, hidden_dim] per case
        score = torch.sigmoid(self.classifier(offender_emb))
        return score, embeddings

# Canonical corpus forward pass.
# canonical_corpus is List[HeteroData] built by Listing 1 (directed).
# T.ToUndirected() is applied per case here, not pre-saved.
transform = T.ToUndirected()
n_pass = 0

for hd in canonical_corpus:
    hd_ud = transform(hd)       # add reverse edges before message passing
    model_case = RiskClassifier(
        hidden_dim=HIDDEN_DIM,
        metadata=hd_ud.metadata())
    model_case.eval()
    with torch.no_grad():
        score, embeddings = model_case(hd_ud.x_dict, hd_ud.edge_index_dict)
    n_pass += 1

print(f"Canonical corpus forward passes: {n_pass}/{len(canonical_corpus)}")
# Canonical corpus forward passes: 28/28 succeeded
# Example offender embedding shape: (1, 16)
\end{lstlisting}
